\documentclass[twocolumn]{aastex631}
\begin{document}
\title{The reionization ionizing-photon-budget shortfall for star-forming galaxies is not robust to systematics at z\~{}6 (beta systematic corner)}
\author{NebulaMind Lab (autonomous pipeline)}
\affiliation{NebulaMind Lab, \url{https://lab.nebulamind.net}}
\begin{abstract}
Reionization ionizing-photon-budget reconciliation at z\~{}6: star-forming galaxies require f\_esc=0.039 (+0.039/-0.020) to close the budget (Madau-Dickinson SFRD, log xi\_ion=25.5+/-0.15, clumping C=2-5, JWST-SFRD tail), versus indirect-proxy-inferred f\_esc=0.050 (+0.076/-0.030) from LzLCS O32/beta calibrations. Median delta(required-inferred)=-0.010 dex-frac (16-84\%: -0.084 to +0.034); 41\% of the systematic MC shows a shortfall. Conclusion: the budget CLOSES within the systematic, and the sign holds under both O32 and beta calibrations. The result is bounded by the xi\_ion x clumping x proxy-calibration systematic, not by statistics. Generated autonomously from public data (jwst) via the NebulaMind Lab runner: a bounded, reproducible, descriptive study.
\end{abstract}
\section{Introduction}
Recent studies have highlighted a potential crisis in the photon budget for reionization, suggesting that current estimates of ionizing photons from star-forming galaxies may not be sufficient to drive the process [Muñoz2024]. This discrepancy has sparked interest in revisiting the calculations and assumptions underlying these estimates. Previous work has emphasized the importance of considering factors such as the clumping factor and the escape fraction of ionizing photons [Davies2021, Park2022], while others have explored the role of galaxy populations at different redshifts [Duncan2015].
\section{Data and method}
In this study, we adopt a literature-anchored budget calculation approach, using the Madau \& Dickinson (2014) analytic fitting function for the cosmic star formation rate density (SFRD). We also rely on published values for xi\_ion and O32/beta f\_esc proxy calibrations from LzLCS [Chisholm+22, Flury+22; Simmonds+24]. Our method focuses on reconciling the ionizing-photon-budget at z\~{}6 by comparing required escape fractions with those inferred from indirect proxies.
\section{Result}
Our analysis reveals that star-forming galaxies require an escape fraction of f\_esc=0.039 (+0.039/-0.020) to close the reionization photon budget, given the Madau-Dickinson SFRD, log xi\_ion=25.5+/-0.15, and a clumping factor C between 2-5. This value is compared to the indirect-proxy-inferred f\_esc=0.050 (+0.076/-0.030) from LzLCS O32/beta calibrations. The median difference between required and inferred escape fractions is -0.010 dex-frac, with a range of -0.084 to +0.034. Notably, 41\% of our systematic Monte Carlo simulations show a shortfall in the photon budget.
\begin{figure}\centering\includegraphics[width=\columnwidth]{result.png}\caption{The reionization ionizing-photon-budget shortfall for star-forming galaxies is not robust to systematics at z\~{}6 (beta systematic corner)}\end{figure}
\section{Caveats}
It is essential to acknowledge the limitations of this study. Our approach relies on automated, single-selection, and uncalibrated measurements, which may introduce biases and uncertainties. The accuracy of our results depends heavily on the assumptions made regarding xi\_ion, clumping factor, and proxy calibrations. Furthermore, our analysis does not account for potential variations in these parameters across different galaxy populations or redshifts. Therefore, while our findings provide valuable insights into the reionization photon budget, they should be interpreted with caution and considered alongside complementary studies that incorporate additional data and more sophisticated modeling techniques.
\section*{References}
\footnotesize
[Muoz2024] Muñoz et al. (2024). Reionization after JWST: a photon budget crisis?. bibcode 2024MNRAS.535L..37M \par [Davies2021] Davies et al. (2021). The Predicament of Absorption-dominated Reionization: Increased Demands on Ionizing Sources. bibcode 2021ApJ...918L..35D \par [Park2022] Park et al. (2022). Calibrating excursion set reionization models to approximately conserve ionizing photons. bibcode 2022MNRAS.517..192P \par [Duncan2015] Duncan et al. (2015). Powering reionization: assessing the galaxy ionizing photon budget at z < 10. bibcode 2015MNRAS.451.2030D \par [Madau2017] Madau (2017). Cosmic Reionization after Planck and before JWST: An Analytic Approach. bibcode 2017ApJ...851...50M

\end{document}
